% Generated draft: requires mathematical and source review.
\documentclass[11pt]{article}
\usepackage[margin=25mm]{geometry}
\usepackage{fontspec}
\setmainfont{Latin Modern Roman}
\usepackage{amsmath,amssymb,mathtools,unicode-math}
\setmathfont{Latin Modern Math}
\usepackage{enumitem,needspace,xurl,hyperref}
\setlength{\parindent}{0pt}
\setlength{\parskip}{5pt}
\setlength{\emergencystretch}{3em}
\newcommand{\N}{\mathbb{N}}
\newcommand{\Z}{\mathbb{Z}}
\newcommand{\Q}{\mathbb{Q}}
\newcommand{\R}{\mathbb{R}}
\newcommand{\C}{\mathbb{C}}
\newcommand{\F}{\mathbb{F}}
\newcommand{\A}{\mathbb{A}}
\newcommand{\PP}{\mathbb{P}}
\newcommand{\E}{\mathbb{E}}
\newcommand{\eps}{\varepsilon}
\newcommand{\dd}{\,\mathrm{d}}
\begin{document}
\section*{UnsolvedMath ID: 800094}
\textbf{Source problem code:} catalogue entry
\subsection{Definitions and mathematical statement}
The supplied source leaves the objects \(S\), \(\widetilde S\), the action \(\varphi\), and the target of the representations \(\rho\) unspecified. The following records the intended characterization problem, subject to these unresolved qualifications.

For vector fields \(U,V\) on \(\widetilde S\) with sufficient regularity to define their Lie bracket, commuting means \([U,V]=0\). Their common zero set is
\[
Z(U,V)=\{p\in\widetilde S:U(p)=0\text{ and }V(p)=0\}.
\]
For a matrix \(A=\begin{pmatrix}a&b\\c&d\end{pmatrix}\), the source defines
\[
A(U,V)=(aU+bV,cU+dV).
\]
It calls a commuting pair \(\rho\)-equivariant when, for every \(p\in\widetilde S\) and \(\gamma\in\pi_1(S)\),
\[
\bigl(U(\varphi(\gamma)(p)),V(\varphi(\gamma)(p))\bigr)
=\rho(\gamma)\bigl(U(p),V(p)\bigr).
\]
As printed, this equality compares tangent vectors based at different points. It requires a specified identification of tangent spaces, or an interpretation involving the differential of \(\varphi(\gamma)\); neither is provided in the supplied record.

Once the missing data and equivariance convention are fixed, let \(\mathcal R\) be the intended class of representations of \(\pi_1(S)\), and let \(\mathcal C_\rho\) be the intended class of admissible \(\rho\)-equivariant commuting vector-field pairs on \(\widetilde S\). The question is to characterize precisely
\[
\left\{\rho\in\mathcal R:
\forall (U,V)\in\mathcal C_\rho,\quad
\exists p\in\widetilde S\text{ such that }U(p)=V(p)=0
\right\}.
\]
The common zero is required on \(\widetilde S\). A fully specified mathematical statement cannot be recovered without the omitted geometric context and equivariance convention.
\subsection{Short English statement}
Characterize those representations for which every admissible equivariant pair of commuting vector fields on \(\widetilde S\) has a common zero. The supplied excerpt omits the geometric hypotheses, representation target, and tangent-space convention needed to make this precise.
\subsection{Sources}
\begin{itemize}
\item Supplied local UnsolvedMath registry record: permanent ID \(800094\), exact source code \texttt{catalogue entry}, title ``Foliations and Diffeomorphism Groups --- Question 35.'' The record attributes the problem to Nariman and Yazdi, with proposed year 2025, and labels it open; that label is supplied metadata, not an independently verified mathematical assessment.
\item Nariman and Yazdi, \emph{Problem List on Foliations and Diffeomorphism Groups} (2025), Question 35. Supplied source URL: \url{https://arxiv.org/abs/2503.18274}. The record reports access on 2026-10-16 and extraction by \texttt{arxiv\_latex}; these are supplied provenance details, not a claim of independent source consultation.
\end{itemize}
\end{document}
